\paragraph{RQ 05 (frequência de deploy e taxa de falha).}
Para o CFR (a), esperamos correlação de Spearman fraca ($|\rho| < 0{,}3$) com a frequência de deploy, em linha com a afirmação do DORA de que velocidade e estabilidade não são um \textit{trade-off} \cite{dora2024report}. A frequência de deploy conta releases, enquanto o CFR (a) conta execuções de CI, e a proporção de execuções com falha depende mais da qualidade dos testes, de testes instáveis e da infraestrutura do CI do que da cadência de releases. Para o CFR (b), esperamos correlação positiva e moderada, contrariando o DORA, mas em parte por construção da métrica: quem publica várias releases por semana quase sempre tem uma nova release em até 7 dias, e basta que ela seja uma correção para que a anterior conte como falha. Esperamos, portanto, que as duas variantes levem a conclusões diferentes sobre o \textit{trade-off}.

\paragraph{RQ 06 (características associadas ao desempenho).}
Esperamos as associações mais fortes com o número de contribuidores e com o tipo do projeto. Repositórios com mais contribuidores recebem mais mudanças e têm mais pessoas para corrigir falhas, o que deve aumentar a frequência de deploy e reduzir o tempo de recuperação. Bibliotecas e \textit{frameworks} publicam versões no próprio GitHub, com \textit{patches} frequentes, enquanto aplicações e serviços costumam entregar por outros canais e publicar releases mais espaçadas, o que deve se refletir na frequência de deploy e no lead time. Para popularidade, linguagem e idade, esperamos diferenças pequenas: todos os repositórios já são populares (mais de 1.000 estrelas), o que restringe a variação, e estrelas medem visibilidade, não processo de entrega. Esperamos tamanhos de efeito pequenos a médios e que poucas comparações continuem significativas após a correção de Holm.

\paragraph{RQ 07 (sensibilidade às definições).}
Esperamos que a classificação geral mude para uma parcela relevante dos repositórios (entre 20\% e 40\% em cada par de combinações), quase sempre para uma categoria vizinha, o que deve resultar em kappa de Cohen ponderado moderado a substancial. Em relação a C1, as combinações C2 e C3 devem deslocar os repositórios para categorias melhores: somar pré-releases ou usar tags aumenta a frequência de deploy, e o lead time (b) tende a ser menor que o (a), que é puxado pelo commit mais antigo de cada release. Esperamos a maior divergência no par C1$\times$C2, que muda as três definições ao mesmo tempo, e a menor no par C2$\times$C3, que compartilha a variante de lead time. Como a categoria geral é a mediana de quatro notas, parte das mudanças em uma única métrica deve ser absorvida, e a variação por métrica deve ser maior que a da classificação geral.

\paragraph{RQ 08 (\textit{rework rate} e evolução temporal).}
Para o \textit{rework rate}, esperamos mediana entre 20\% e 40\% de releases corretivas, porque, em projetos que seguem o Versionamento Semântico, versões de \textit{patch} são comuns. Esperamos que ele seja maior que o CFR (b) e fortemente correlacionado com ele: o CFR (b) atribui a falha à release que precisou de correção e só considera correções em até 7 dias, enquanto o \textit{rework rate} conta as próprias releases corretivas, inclusive as publicadas depois desse prazo e as que corrigem outra correção. Para a evolução temporal, esperamos que praticamente nenhum repositório apresente tendência significativa em qualquer direção. Um ano é pouco tempo para mudanças de processo e, com apenas quatro trimestres por repositório, o teste de Mann-Kendall bilateral não consegue atingir $p < 0{,}05$ (o menor p-valor possível é cerca de 0,09), de modo que a ausência de tendência refletirá também a falta de poder do teste.
